\section{Implications of Advances in Large Language Models}\label{sec:llm-advances}

Advances in large language models warrant reassessing design choices motivated by the limitations of earlier models. For example, Claude Fable 5.1 and GPT-6 Astra support context windows of one million and 1.05 million tokens, respectively~\cite{anthropic2026models,openai2026models}. Larger context windows may reduce the need to split documentation or source code into smaller inputs. Improvements in task-specific reliability may also reduce the need for corrective prompting and repair. Some supporting mechanisms adopted for earlier models may therefore become unnecessary. 

The appropriate simplifications depend on both the task and the codebase. Large codebases can substantially exceed even million-token context windows. Google's monolithic repository, for example, contains billions of lines of code~\cite{potvin2016google,christopher2025migration}. Selecting task-relevant information remains important at this scale. Even when an input fits within the window, its effective use requires separate evaluation. NoLiMa minimizes lexical overlap between questions and relevant passages to test retrieval through semantic associations. Retrieval performance declined with increasing context length in its evaluated models~\cite{modarressi2025nolima}, motivating task-specific evaluation of long-context models.

The following discussion considers which parts of each system could be simplified and which retain value through executable checks or reusable program representations. It also considers whether these components can reduce total computation at comparable task accuracy. The proposed changes and their benefits are prospective, distinct from the implementations and measurements reported in the preceding chapters.

\subsection{CompSuite: Reproducible Evaluation}
\toolcompsuite provides reproducible library-upgrade incompatibility cases and an execution tool. As a dataset and evaluation infrastructure, it does not rely on an LLM inference pipeline whose components would need to be simplified as models improve.

Stronger models may reduce the guidance and repair iterations needed to resolve the incompatibilities recorded in \toolcompsuite. Its fixed client-library versions, reproducible environments, and developer-written tests would still provide a common basis for evaluating these workflows. If future systems solve every recorded case, the benchmark would become less informative for distinguishing repair effectiveness, but could still support comparisons of model calls, token consumption, and elapsed time required to achieve the same test outcomes. These outcomes establish success on the available tests rather than correctness across all possible executions. More broadly, even when both library and client code are generated entirely by AI, their independent evolution can still introduce incompatible interfaces or behaviors. The recorded cases therefore remain useful as concrete failure patterns for compatibility analysis and regression testing.

\subsection{MPChecker: Simplifying Constraint Extraction}
The documentation processing in \toolchecker offers concrete opportunities for simplification. The 1,500-word chunks in \sref{sec:llm-mpchecker} were selected for the evaluated GPT-4 configuration. Larger context windows could allow complete API descriptions to be processed together, reducing fragmentation and repeated prompt instructions. More reliable constraint extraction could also support fewer examples and shorter prompts. These changes may reduce model calls and repeated input, although processing unrelated documentation can offset those savings. Input selection should therefore depend on the information needed for extraction rather than the maximum available context window.

Better extraction also warrants reassessing fuzzy processing. Similarity-based tolerance of incorrect parameter names or operators could be reduced or removed if a newer model produces sufficiently accurate constraints. Predicates for documentation terms such as ``ignore'' and ``exist'' serve a different purpose by representing relationships between parameter use and execution paths. Those semantics still need to be represented, but their current fuzzy implementation is not necessarily the only suitable approach. The ablation in \sref{sec:eval-mpchecker} removes both fuzzy words and fuzzy constraint logic, so it does not establish which parts could be removed as extraction improves.

Explicit constraints, symbolic execution, and SMT checking retain value even as extraction becomes more accurate. Within the scope of the represented constraints and analyzed paths, these mechanisms provide inspectable evidence of whether program behavior satisfies or violates documented requirements. Such evidence supports review by both developers and AI agents in human--AI collaboration and fully automated software engineering workflows. During maintenance of large systems, these checks can help identify overlooked violations of textual constraints and direct further investigation to the relevant code. By executing explicit checks rather than repeatedly asking an LLM to infer whether code satisfies each requirement, automated software engineering systems may reduce their reliance on model reasoning and the associated inference costs. 

\subsection{DataLoc: Fewer Repairs with Reusable Program Facts}
For \tooldataloc, more reliable Datalog generation could reduce synthesis and repair iterations. Larger contexts could accommodate more relevant predicate definitions and source snippets during query interpretation, potentially reducing errors caused by missing context. The mutation diagnostics already operate on zero-cardinality intermediate relations under a bounded budget. With a newer model, a smaller probe budget or removal of repair heuristics that provide little additional benefit could simplify this stage and reduce diagnostic executions and subsequent model calls. These changes should be evaluated on both matching and empty-ground-truth queries, since syntactically valid rules may still misrepresent the user's intent.

Program facts and the Datalog engine have a role beyond compensating for generation errors. For queries expressible over the extracted facts, the engine evaluates structural relationships without requiring the LLM to reconstruct them from repository text for each query. The same facts can support multiple queries over a repository version. Parser validation also provides an explicit check of generated rule syntax. These components provide reusable computation and executable checks, although their costs must be included when assessing the overall benefit. An empty result still requires careful interpretation because incomplete facts or an incorrect rule may omit intended matches. If mutation probes are retained, their relaxed results should remain diagnostic feedback rather than become final answers.

Source-level candidate verification is another component to reassess because it incurs additional model inference. Its contribution could be measured by comparing the full pipeline with a variant that omits this step under the same generation model. If verification contributes little to accuracy for a newer model, removing it could reduce per-query cost. If it continues to identify mismatches between generated rules and user intent, retaining it may justify that cost. Improved generation alone does not determine which choice is preferable.

\subsection{Computational Efficiency}
A potential computational benefit of a neurosymbolic architecture comes from using explicit program representations to support repeated operations. The LLM interprets natural language and formulates constraints or queries, while symbolic tools perform structural traversal, relational filtering, and constraint checking. This division can reduce repeated model inference over source code. More reliable generation can complement this design by reducing the repair work needed before a query or constraint can be executed.

The benefit depends on the workload. In \tooldataloc, extracting facts once allows that cost to be spread across multiple queries over the same repository version. Reusing code constraints in \toolchecker could offer a similar benefit when checking revised documentation, although this extension remains to be implemented and evaluated. For small codebases or one-off tasks, the initial extraction and analysis costs may outweigh the savings from symbolic execution. Frequent repository changes can also increase maintenance costs. A simpler model-only workflow may therefore be preferable in some settings, even if reuse benefits the hybrid architecture in others.

The measurements in \sref{sec:cost-dataloc} illustrate the trade-off in the evaluated setting. On \dataset, \tooldataloc with Claude-3.5 averages 39.3 seconds and 16.2k tokens per query, compared with 393.0k tokens for LocAgent with Claude-3.5. This reduction in token use accompanies stronger localization results on KA-LCL. Agentless with Claude-3.5 uses 15.2k tokens, showing that lower token consumption alone does not establish an advantage. It must be considered alongside task accuracy. These measurements motivate further comparison, but do not establish computational savings with newer models or quantify total computation from token counts alone.

Further evaluation should compare direct long-context inference, simplified hybrid variants, and the full architecture using the same model and tasks. Measuring accuracy alongside model calls, input and output tokens, elapsed time, and symbolic analysis and execution costs would reveal the trade-offs. Initial extraction, repository updates, and comparable caching opportunities should be included for all approaches. Varying repository size and the number of queries would help identify when reuse offsets its initial cost. This evaluation would establish which supporting mechanisms can be removed and when retaining symbolic components reduces total computation at comparable task accuracy.
